\documentclass[10pt,conference,letterpaper]{IEEEtran}
\usepackage[T1]{fontenc}
\usepackage{amsmath,amssymb}
\usepackage{url}
\usepackage[hidelinks]{hyperref}
\usepackage[nospread,noshrink]{flushend}

\newtheorem{proposition}{Proposition}
\newcommand{\codefile}[2]{\href{https://github.com/DongXinfzu/LEAP/blob/main/#1}{\texttt{#2}}}

\hypersetup{
  pdftitle={LEAP: Supplementary Proofs},
  pdfauthor={Xin Dong, Yusen Liu, Hua Lu, Huan Li, Dalin Zhang, Zhichen Lai, Xing Chen},
  pdfsubject={Repository appendix: geometric and DTW lower-bound proofs},
  pdfkeywords={LEAP, trajectory similarity search, lower bounds, DTW}
}

\title{LEAP: Supplementary Proofs}
\author{\IEEEauthorblockN{Xin Dong\textsuperscript{1,*}, Yusen Liu\textsuperscript{2,*},
Hua Lu\textsuperscript{3}, Huan Li\textsuperscript{4}, Dalin Zhang\textsuperscript{5},
Zhichen Lai\textsuperscript{1,\textdagger}, and Xing Chen\textsuperscript{1}}
\IEEEauthorblockA{\textsuperscript{1}Fuzhou University;
\textsuperscript{2}University of Technology Sydney;
\textsuperscript{3}Aalborg University\\
\textsuperscript{4}Zhejiang University;
\textsuperscript{5}Hangzhou Dianzi University\\
\textsuperscript{*}Equal contribution.
\textsuperscript{\textdagger}Corresponding author.}}

\begin{document}
\maketitle

\begin{abstract}
This repository companion collects the complete proofs of LEAP's cross-field
lower bound, critical-position proximity property, weighted DTW block bound,
and DTW signal lower bound. The four appendices correspond to Appendices A--D
of the repository README. All statements concern real arithmetic; the
implementation links identify the corresponding computations.
\end{abstract}

\section*{Notation}
For a trajectory $T$, let $V_T$ be its sampled positions and $E_T$ the union
of the line segments joining consecutive positions. The distance from a point
$x$ to a set $S$ is
$\mathrm{dist}(x,S)=\min_{z\in S}\lVert x-z\rVert_2$.
For probe $p_j^{\mathrm{pr}}$, the nearest-position and segment fields are
\begin{equation}
\begin{aligned}
a_j(T)&=\mathrm{dist}(p_j^{\mathrm{pr}},V_T),\\
b_j(T)&=\mathrm{dist}(p_j^{\mathrm{pr}},E_T).
\end{aligned}
\end{equation}
The directed and symmetric vertex-to-segment Hausdorff distances are
\begin{equation}
\begin{aligned}
h(T_A,T_B)&=\max_{p\in V_{T_A}}\mathrm{dist}(p,E_{T_B}),\\
D_H(T_A,T_B)&=\max\{h(T_A,T_B),h(T_B,T_A)\}.
\end{aligned}
\end{equation}
The cross-field lower bound combines all probes and both directions:
\begin{equation}
\begin{aligned}
L_H(T_A,T_B)=\max\bigl\{0,\;&\max_j[b_j(T_B)-a_j(T_A)],\\
&\max_j[b_j(T_A)-a_j(T_B)]\bigr\}.
\end{aligned}
\end{equation}

\appendices
\section{Validity of the Cross-Field Bound}
\label{app:validity}

\begin{proposition}
For any probe positions,
$0\leq L_H(T_A,T_B)\leq D_H(T_A,T_B)$ in real arithmetic.
Validity does not depend on learning.
\end{proposition}

\begin{IEEEproof}
Fix $p_j^{\mathrm{pr}}$ and choose its nearest position $p\in V_{T_A}$.
Apply the Euclidean triangle inequality through $p$ to its nearest point on
$E_{T_B}$ to obtain
\begin{equation}
\begin{aligned}
b_j(T_B)&=\mathrm{dist}(p_j^{\mathrm{pr}},E_{T_B})\\
&\leq\lVert p_j^{\mathrm{pr}}-p\rVert_2+\mathrm{dist}(p,E_{T_B})\\
&\leq a_j(T_A)+h(T_A,T_B).
\end{aligned}
\end{equation}
Hence $b_j(T_B)-a_j(T_A)\leq h(T_A,T_B)$. Interchanging $T_A$ and $T_B$
proves the other direction. Taking the maximum over probes and both directions,
with zero included, yields the result.
\end{IEEEproof}

\section{Critical-Position Proximity}
\label{app:proximity}

A \emph{critical position} $p^*$ is a sampled position that attains
$D_H(T_A,T_B)$ in either direction. For probe set $\Theta$, let
$\delta_\Theta(p^*)=\min_j\lVert p_j^{\mathrm{pr}}-p^*\rVert_2$ be its
distance to the nearest probe.

\begin{proposition}
For any critical position $p^*$,
\begin{equation}
0\leq D_H(T_A,T_B)-L_H(T_A,T_B)\leq 2\delta_\Theta(p^*).
\end{equation}
\end{proposition}

\begin{IEEEproof}
Suppose $p^*\in V_{T_A}$ and
$\mathrm{dist}(p^*,E_{T_B})=D_H(T_A,T_B)$. Choose the nearest probe
$p_j^{\mathrm{pr}}$, at distance $\delta=\delta_\Theta(p^*)$.
By the same triangle-inequality argument,
$b_j(T_B)\geq D_H(T_A,T_B)-\delta$. Since $p^*$ is a position of $T_A$,
$a_j(T_A)\leq\delta$. Therefore
\begin{equation}
\begin{aligned}
L_H(T_A,T_B)&\geq b_j(T_B)-a_j(T_A)\\
&\geq D_H(T_A,T_B)-2\delta.
\end{aligned}
\end{equation}
The case $p^*\in V_{T_B}$ follows by symmetry, and validity gives nonnegativity.
\end{IEEEproof}

\section{Weighted DTW Block Bound}
\label{app:block}

This bound is a check before original-distance refinement, not the signal
bound used for candidate ranking. It follows the paper's DTW subsection
under Filtering and Exact Search.

\subsection{Blocks and Costs}
Let $q=(p_1^q,\ldots,p_n^q)$ and $T=(p_1^T,\ldots,p_m^T)$ be nonempty
trajectories. Their DTW distance $D_W(q,T)$ is the minimum sum of Euclidean
point-to-point distances along an alignment from $(1,1)$ to $(n,m)$, using
steps $(1,0)$, $(0,1)$, or $(1,1)$. Thus every position of each trajectory is
matched at least once, and positions are visited in order.

We split each trajectory into four consecutive blocks with nearly equal
position counts. For a trajectory of length $n\geq4$, block $s$ contains
positions $\lceil(s-1)n/4\rceil+1$ through $\lceil sn/4\rceil$, for
$s\in\{1,2,3,4\}$. If $n<4$, each position forms one block and the remaining
slots are empty. Query block $s$ contains $n_s$ positions, and candidate block
$t$ contains $m_t$ positions.

For each nonempty block and fixed filtering probe $p_{W,j}^{\mathrm{pr}}$,
we record the minimum and maximum distance from its positions to that probe.
The intervals are $[l^q_{sj},h^q_{sj}]$ for the query and $[l^T_{tj},h^T_{tj}]$
for the candidate. Rows of the following matrix index query blocks; columns
index candidate blocks:
\begin{equation}
G_{st}=\max_j\max\{0,l^q_{sj}-h^T_{tj},l^T_{tj}-h^q_{sj}\}.
\end{equation}
Empty blocks have zero counts and zero intervals, and we explicitly set
$G_{st}=0$ whenever either block is empty. For positions $x$ and $y$ in
nonempty blocks $s$ and $t$, respectively, the reverse triangle inequality gives
\begin{equation}
\begin{aligned}
\lVert x-y\rVert_2
&\geq\bigl|\lVert x-p_{W,j}^{\mathrm{pr}}\rVert_2
          -\lVert y-p_{W,j}^{\mathrm{pr}}\rVert_2\bigr|\\
&\geq\max\{0,l^q_{sj}-h^T_{tj},l^T_{tj}-h^q_{sj}\}.
\end{aligned}
\end{equation}
Taking the maximum over probes shows that $G_{st}$ lower-bounds every
point-pair distance in those blocks. For a nonnegative block-cost matrix $A$,
write $\mathrm{DTW}(A)$ for the minimum summed cost along a path from its
upper-left to lower-right cell, with the same three allowed steps.
The coarse bound is $\mathrm{DTW}(G)$.

\subsection{Weighted Bounds}
For each query block, its smallest cost over candidate blocks lower-bounds
the matching cost of each query position in that block. Weighting this
minimum by the block's position count accounts for those positions. We
subtract this minimum before processing candidate blocks so that their
contributions do not count the same part of a cost twice:
\begin{equation}
\begin{aligned}
u_s&=\min_t G_{st},\qquad v_t=\min_s(G_{st}-u_s),\\
R_{st}&=G_{st}-u_s-v_t\geq0,\\
L_{\mathrm{row}}&=\sum_s n_su_s+\sum_t m_tv_t+\mathrm{DTW}(R).
\end{aligned}
\end{equation}
Here $u_s$ is the query-block contribution per position, $v_t$ is the
candidate-block contribution after subtraction, and $R$ contains the
remaining costs. Its DTW term retains the order of blocks. The subscript
\texttt{row} means that query blocks (matrix rows) are processed first.

Processing candidate blocks (matrix columns) first gives a second bound:
\begin{equation}
\begin{aligned}
v'_t&=\min_s G_{st},\qquad u'_s=\min_t(G_{st}-v'_t),\\
R'_{st}&=G_{st}-u'_s-v'_t\geq0,\\
L_{\mathrm{col}}&=\sum_s n_su'_s+\sum_t m_tv'_t+\mathrm{DTW}(R').
\end{aligned}
\end{equation}
We keep the strongest of the coarse and two weighted bounds:
\begin{equation}
L_{\mathrm{block}}
=\max\{\mathrm{DTW}(G),L_{\mathrm{row}},L_{\mathrm{col}}\}.
\end{equation}

\subsection{Validity}
\begin{proposition}
In real arithmetic, $0\leq L_{\mathrm{block}}(q,T)\leq D_W(q,T)$.
\end{proposition}

\begin{IEEEproof}
Fix any original alignment $\pi$. For a matched index pair $(i,r)\in\pi$,
let $s(i)$ and $t(r)$ be the blocks containing positions $p_i^q$ and $p_r^T$.
The block-cost inequality gives
\begin{equation}
\begin{aligned}
\sum_{(i,r)\in\pi}\lVert p_i^q-p_r^T\rVert_2
&\geq\sum_{(i,r)\in\pi}G_{s(i),t(r)}\\
&=\sum_{(i,r)\in\pi}
\bigl(u_{s(i)}+v_{t(r)}+R_{s(i),t(r)}\bigr).
\end{aligned}
\end{equation}
Every query position is matched at least once. Since $u_s\geq0$, the sum
of query contributions is at least $\sum_s n_su_s$. Similarly, $v_t\geq0$
gives a candidate contribution of at least $\sum_t m_tv_t$. These sums count
different portions of each matched-pair cost, not two copies of that cost.

Map the alignment to its block pairs and collapse consecutive repetitions
of the same pair. Consecutive blocks preserve the original order, so this
produces a valid block path. Removing repetitions cannot increase the sum
of the nonnegative remaining costs $R_{st}$. Its cost is at least
$\mathrm{DTW}(R)$, the minimum over all such paths. Consequently,
\begin{equation}
\begin{aligned}
\sum_{(i,r)\in\pi}\lVert p_i^q-p_r^T\rVert_2
&\geq\sum_s n_su_s+\sum_t m_tv_t\\
&\quad+\mathrm{DTW}(R)=L_{\mathrm{row}}.
\end{aligned}
\end{equation}
The inequality holds for every original alignment, hence
$L_{\mathrm{row}}\leq D_W(q,T)$. Reversing the subtraction order proves
$L_{\mathrm{col}}\leq D_W(q,T)$. Applying the same block-path argument
directly to $G$ proves $\mathrm{DTW}(G)\leq D_W(q,T)$. Taking their maximum
preserves the lower-bound inequality.

For trajectories with fewer than four positions, the mapped path ends at
the last nonempty block pair. Extend it to the final padded cell through
zero-cost cells. On every padded cell, $G_{st}=0$ and the nonnegative
decomposition forces $u_s=v_t=R_{st}=0$; the column-first case is identical.
Thus padding adds no cost, and the argument still applies.
\end{IEEEproof}

Filtering can therefore stop early when
$L_{\mathrm{block}}(q,T)>\widehat r_{W,q}$. Equality does not trigger
rejection. Remaining candidates proceed to original DTW refinement with
its running-bound early stopping.

\paragraph{Implementation}
The block-filter implementation provides
\codefile{01_LEAP/leap/dtw/dtw_probe_intervals.py}{counts} and
\codefile{01_LEAP/leap/dtw/dtw_probe_intervals.py}{stats\_torch}
for the block partition,
\codefile{01_LEAP/leap/dtw/dtw_probe_intervals.py}{bound\_torch}
for the reference calculation, and
\codefile{01_LEAP/leap/dtw/dtw_probe_intervals.py}{pair\_kernel}
for the GPU calculation. The existing implementation also applies a
dtype-dependent downward numerical guard; the proof above concerns the
mathematical bound in real arithmetic.

\section{DTW Signal Lower Bound}
\label{app:signal}

This is the signal bound used to rank DTW candidates
(\texttt{eq:dtw-signal}), not the block bound in Appendix~\ref{app:block}
or the learned summary retrieval score.

\subsection{Definition}
Let $A=(x_1,\ldots,x_n)$ and $B=(y_1,\ldots,y_m)$ be nonempty trajectories,
and let $c_{W,1},\ldots,c_{W,M_s}$ be a shared set of $M_s\geq1$ probes.
For a vertex $x$, its signal at probe $j$ is
$\sigma_j(x)=\lVert x-c_{W,j}\rVert_2$. Define the signal matching cost by
\begin{equation}
\ell_{\mathrm{sig}}(x,y)
=\max_{1\leq j\leq M_s}|\sigma_j(x)-\sigma_j(y)|.
\end{equation}
Let $\Pi_{n,m}$ contain the monotone paths from $(1,1)$ to $(n,m)$ with
steps $(1,0)$, $(0,1)$, or $(1,1)$. Both quantities use this same path set,
retaining every vertex in its original order:
\begin{equation}
\begin{aligned}
L_W^{\mathrm{sig}}(A,B)
&=\min_{\pi\in\Pi_{n,m}}\sum_{(i,l)\in\pi}
\ell_{\mathrm{sig}}(x_i,y_l),\\
D_W(A,B)
&=\min_{\pi\in\Pi_{n,m}}\sum_{(i,l)\in\pi}
\lVert x_i-y_l\rVert_2.
\end{aligned}
\end{equation}
Here DTW sums Euclidean distances, without squaring or length normalization.
Earlier manuscript notation $D_W^{\mathrm{sig}}$ denotes the same signal
quantity as $L_W^{\mathrm{sig}}$ above.

\subsection{Validity}
\begin{proposition}
For any shared probe positions, in real arithmetic,
$0\leq L_W^{\mathrm{sig}}(A,B)\leq D_W(A,B)$.
\end{proposition}

\begin{IEEEproof}
The Euclidean triangle inequality gives, for every probe $j$,
\begin{equation}
\begin{aligned}
\sigma_j(x)&\leq\lVert x-y\rVert_2+\sigma_j(y),\\
\sigma_j(y)&\leq\lVert x-y\rVert_2+\sigma_j(x).
\end{aligned}
\end{equation}
Combining these two inequalities yields
$|\sigma_j(x)-\sigma_j(y)|\leq\lVert x-y\rVert_2$. The right-hand side
does not depend on $j$, so taking the maximum over probes preserves it:
$0\leq\ell_{\mathrm{sig}}(x,y)\leq\lVert x-y\rVert_2$.

Now choose an original-DTW minimizing path $\pi_W\in\Pi_{n,m}$, which
exists because the trajectories and path set are finite. The signal minimum
cannot exceed the signal cost along this particular path. Summing the
point-pair inequalities along it gives
\begin{equation}
\begin{aligned}
0\leq L_W^{\mathrm{sig}}(A,B)
&\leq\sum_{(i,l)\in\pi_W}\ell_{\mathrm{sig}}(x_i,y_l)\\
&\leq\sum_{(i,l)\in\pi_W}\lVert x_i-y_l\rVert_2\\
&=D_W(A,B).
\end{aligned}
\end{equation}
The signal-minimizing path need not equal $\pi_W$. No interchange of the
minimum, sum, or maximum is used. Repeated vertices, ties, and single-vertex
trajectories do not change the argument. Learning can change the bound's
tightness, but validity holds for any shared probe positions.
\end{IEEEproof}

\paragraph{Implementation}
The signal encoder
\codefile{01_LEAP/leap/dtw/dtw_signal.py}{encode}
computes vertex-to-probe distances. The
\codefile{01_LEAP/leap/dtw/dtw_signal.py}{signal\_kernel}
takes the maximum absolute signal difference for each matched pair and
applies the three-predecessor DTW recurrence. The
\codefile{01_LEAP/leap/dtw/dtw_signal_bucket.py}{pairs}
function buckets trajectory lengths without dropping vertices or changing
the recurrence. The proof concerns exact real arithmetic, not a guarantee that
floating-point roundoff preserves the inequality without a numerical guard.
It does not make the separate learned summary score a lower bound.

\section*{Repository}
The README proofs and corresponding implementation are available at
\url{https://github.com/DongXinfzu/LEAP}.

\end{document}
